\begin{lemma}[Missing and stored colors are monotone through one series iteration]
Let \(\mathcal T\) be a certified Tashkinov series with \(r\) noninitial
iterations, and let \(q<r\).  Put
\[
 n=q+1,\quad T_n=\mathcal T.\mathrm{tree}(q+1),\quad
 T_{n+1}=\mathcal T.\mathrm{tree}(q+2),
\]
\[
 \varphi_{n-1}=\mathcal T.\mathrm{coloring}(q),\quad
 \varphi_n=\mathcal T.\mathrm{coloring}(q+1),\quad
 D_{n-1}=D_q,\quad D_n=D_{q+1}.
\]
Then both inclusions
\[
 \overline\varphi_{n-1}(T_n)\cup D_{n-1}
 \subseteq \overline\varphi_n(T_n)\cup D_n
 \subseteq \overline\varphi_n(T_{n+1})\cup D_n
\]
hold.
\end{lemma}
\begin{proof}
Write \(S_{\leq a}=\bigcup_{0\leq h\leq a}S_h\), using the exact
stored sets carried by the series.  We first record two set identities.
By \noderef{TashkinovStoredColors}, for a positive index \(a\),
\(D_a=S_{\leq a}\setminus\overline\varphi_a(T_a)\), whereas
\(D_0=\varnothing\).  Every certified noninitial constructor in
\noderef{TashkinovSeriesStep} makes the old tree an ordered initial
segment of the new tree.  At index zero, the initial constructor makes
\(T_1\) a \noderef{TashkinovClosure} of the seed \(T_0\).  Expanding
\noderef{TashkinovAugmentationSequence}, every step in this closure is a
one-edge \noderef{TashkinovAugmentation} and hence an ordered
initial-segment extension; an
induction on the reflexive-transitive closure therefore also makes
\(T_0\) an initial segment of \(T_1\).  By
\noderef{TreeSequenceInitialSegment}, all vertices of an earlier tree
occur in the later tree.  The union definition
\noderef{TreeMissingColors} consequently gives
\[
 \overline\psi(T_a)\subseteq\overline\psi(T_{a+1})                 \tag{1}
\]
for every coloring \(\psi\) and every consecutive pair of certified
trees.

Apply the elementary identity
\(A\cup(B\setminus C)=A\cup B\) when \(C\subseteq A\).  With
\(A=\overline\varphi_{n-1}(T_n)\), (1) supplies
\(\overline\varphi_{n-1}(T_{n-1})\subseteq A\), and hence
\[
 \overline\varphi_{n-1}(T_n)\cup D_{n-1}
 =\overline\varphi_{n-1}(T_n)\cup S_{\leq n-1}.                  \tag{2}
\]
When \(n=1\), the same equality follows from the initial-constructor
equalities \(D_0=S_0=\varnothing\) in
\noderef{TashkinovSeriesStep}.  Since the set subtracted in the definition
of \(D_n\) is already \(\overline\varphi_n(T_n)\), we likewise have
\[
 \overline\varphi_n(T_n)\cup D_n
 =\overline\varphi_n(T_n)\cup S_{\leq n}.                        \tag{3}
\]

It remains to compare the missing colors on the fixed tree \(T_n\).
Use the certificate field of \noderef{TashkinovSeries} at the positive
index \(q+1\), and distinguish its three possible constructors from
\noderef{TashkinovSeriesStep}.

In an RE, the constructor states
\(\varphi_n=\varphi_{n-1}\), so every old missing color is still missing.
In an SE, \(\varphi_n\) is the selected coloring, which is
\((T_n,D_{n-1},\varphi_{n-1})\)-stable.  The vertexwise equality in
\noderef{StableColoring} therefore gives
\(\overline\varphi_n(T_n)=\overline\varphi_{n-1}(T_n)\).

In a PE, call the selected stable coloring \(\pi\), the second stable
coloring \(\rho\), the supporting vertex \(v_n\), and the selected colors
\(\gamma_n,\delta_n\).  The two stable-coloring conditions give, at each
vertex of \(T_n\), first
\(\overline\pi(x)=\overline\varphi_{n-1}(x)\) and then
\(\overline\rho(x)=\overline\pi(x)\).  The PE certificate also gives a
\noderef{TreeExit} for the \((\gamma_n,\delta_n)\)-component through
\(v_n\), with no other tree vertex on its exit path by
\noderef{TreeExitPath}.  Here \noderef{ParallelExtensionColor} says that
\(\gamma_n\) is missing at \(v_n\), while the selected connecting edge is
incident there and has color \(\delta_n\) under \(\rho\).  Thus
\(\gamma_n\ne\delta_n\), and \(v_n\) has degree one in the two-color
incidence graph.  The outside end of the exit path also has degree one:
one of the two colors is missing there, while the nontrivial path supplies
an incident edge of the other color.  Properness and
\noderef{TwoColorIncidenceDegreeAtMostTwo} bound every intermediate degree
by two.  A simple path between two degree-one vertices in a connected
graph of maximum degree two exhausts their component: otherwise the first
edge from the path to a node outside it would give degree at least three at
an internal node, or degree at least two at an endpoint.  Hence the full
component meets \(V(T_n)\) only at \(v_n\).

The definition \noderef{TwoColorComponentSwap} now shows that the PE
switch changes no tree-vertex missing-color set except at \(v_n\), where
it replaces the missing color \(\gamma_n\) by \(\delta_n\).  Consequently
every member of \(\overline\varphi_{n-1}(T_n)\) remains in
\(\overline\varphi_n(T_n)\), except possibly \(\gamma_n\); in that
exceptional case \(\gamma_n\in S_n=\{\delta_n,\gamma_n\}\).  The RE and
SE cases satisfy the same conclusion without an exception.  Therefore
\[
 \overline\varphi_{n-1}(T_n)
 \subseteq\overline\varphi_n(T_n)\cup S_n.                       \tag{4}
\]

Combining (2), (4), and
\(S_{\leq n-1}\cup S_n=S_{\leq n}\), and then using (3), proves the
first asserted inclusion.  Finally, (1), now for \(T_n\subseteq T_{n+1}\)
under the fixed coloring \(\varphi_n\), gives
\(\overline\varphi_n(T_n)\subseteq\overline\varphi_n(T_{n+1})\).
Adjoining the unchanged set \(D_n\) to both sides proves the second
inclusion.
\end{proof}
